\DeclareRobustCommand{\code}[1]{\textcolor{BookCodeInline}{\texttt{\nolinkurl{#1}}}}
\DeclareRobustCommand{\asmcode}[1]{\textcolor{BookCodeInline}{\texttt{\nolinkurl{#1}}}}
\DeclareRobustCommand{\filepath}[1]{\textcolor{BookCodeInline}{\texttt{\nolinkurl{#1}}}}
\DeclareRobustCommand{\term}[2]{\textbf{#1}\,（#2）}
\providecommand{\cpp}{\texorpdfstring{\mbox{\sffamily C\kern-0.06em+\kern-0.06em+}}{C++}}

\newcolumntype{L}[1]{>{\raggedright\arraybackslash}p{#1}}
\newcolumntype{F}[1]{>{\raggedright\arraybackslash}p{\dimexpr0.94\dimexpr#1\relax-0.55pt\relax}}
\setlength{\arrayrulewidth}{0.65pt}
\setlength{\heavyrulewidth}{1.00pt}
\setlength{\lightrulewidth}{0.65pt}
\setlength{\cmidrulewidth}{0.55pt}
\arrayrulecolor{BookTableRule}
\renewcommand{\arraystretch}{1.18}
\newcommand{\tablerowrule}{\hline}

% 文件系统章采用与硬件书一致的四层版式：少量大区段、块级主题标题、
% 更轻的子主题标题，以及在手机缩放下仍能分辨列边界的表格线。
\newenvironment{filesystemlayout}
{%
  \begingroup
  \setlength{\arrayrulewidth}{0.65pt}%
  \setlength{\heavyrulewidth}{1.00pt}%
  \setlength{\lightrulewidth}{0.65pt}%
  \setlength{\cmidrulewidth}{0.55pt}%
  \arrayrulecolor{BookTableRule}%
  \renewcommand{\arraystretch}{1.18}%
  \tikzset{
    os block/.append style={rounded corners=2pt,line width=0.42pt,inner xsep=5pt,inner ysep=4pt},
    os state/.style={os block,draw=BookTeal,fill=BookCodeBg},
    os compute/.style={os block,draw=BookAccent,fill=BookCodeBg},
    os memory/.style={os block,draw=BookAmber,fill=white},
    os control/.style={os block,draw=BookRed,fill=white}
  }%
}
{%
  \endgroup
}

\newcommand{\fsmajorsection}[1]{%
  \section{#1}%
}

\newcommand{\fsdivision}[1]{%
  \section{#1}%
}

\newcommand{\fstopic}[1]{%
  \Needspace{4\baselineskip}%
  \par\addvspace{0.72em}%
  {\noindent\normalsize\bfseries\color{BookNavy}#1\par}%
  \nobreak\smallskip
}

\newcommand{\fssubtopic}[1]{%
  \Needspace{3\baselineskip}%
  \par\addvspace{0.48em}%
  {\noindent\normalsize\bfseries\color{BookAccent}#1\par}%
  \nobreak\smallskip
}

\newenvironment{fspdfdiagram}[1][]
{\par\vspace{1.12em plus 0.22em}\goodbreak\noindent\begin{center}\begin{tikzpicture}[os diagram,scale=0.94,transform shape,#1]}
{\end{tikzpicture}\end{center}\vspace{0.68em plus 0.16em}}

\newcommand{\fsdiagramnote}[1]{%
  \par\goodbreak
  {\noindent\footnotesize\color{BookMuted}图示：#1\par}%
  \vspace{0.72em plus 0.14em}%
}

\lstdefinestyle{oscode}{
  basicstyle=\ttfamily\small,
  backgroundcolor=\color{BookCodeBg},
  frame=single,
  rulecolor=\color{BookRule},
  columns=fullflexible,
  keepspaces=true,
  breaklines=true,
  showstringspaces=false,
  xleftmargin=0.8em,
  xrightmargin=0.4em,
  aboveskip=0.75em,
  belowskip=0.75em
}
\lstset{style=oscode}

\newtcolorbox{keyidea}{
  enhanced,
  breakable,
  sharp corners,
  boxrule=0pt,
  leftrule=0.5pt,
  colback=BookCodeBg,
  colframe=BookAccent,
  title=核心思想,
  fonttitle=\bfseries,
  coltitle=BookNavy,
  before skip=0.75em,
  after skip=0.75em
}

\newtcolorbox{warningbox}{
  enhanced,
  breakable,
  sharp corners,
  boxrule=0pt,
  leftrule=0.5pt,
  colback=white,
  colframe=BookRed,
  title=常见误区,
  fonttitle=\bfseries,
  coltitle=white,
  before skip=0.75em,
  after skip=0.75em
}

\newtcolorbox{rubricbox}{
  enhanced,
  breakable,
  sharp corners,
  boxrule=0pt,
  leftrule=0.5pt,
  colback=white,
  colframe=BookAmber,
  title=验收标准,
  fonttitle=\bfseries,
  coltitle=white,
  before skip=0.75em,
  after skip=0.75em
}

\newtcolorbox{statebox}{
  enhanced,
  breakable,
  sharp corners,
  boxrule=0pt,
  leftrule=0.6pt,
  colback=BookBlueTint,
  colframe=BookAccent,
  left=1.2em,
  right=1.0em,
  top=0.65em,
  bottom=0.65em,
  before skip=0.60em,
  after skip=0.62em
}

\newtcolorbox{problemframe}{
  enhanced,
  breakable,
  sharp corners,
  boxrule=0pt,
  leftrule=0.6pt,
  colback=BookAmberTint,
  colframe=BookAmber,
  left=1.2em,
  right=1.0em,
  top=0.65em,
  bottom=0.65em,
  before skip=0.60em,
  after skip=0.62em
}

% 章末练习采用教材式编号题目，不使用表格承载题干和答案。
\newcounter{osexercise}[chapter]
\newcounter{osanswer}[chapter]
\renewcommand{\theosexercise}{\thechapter.\arabic{osexercise}}
\renewcommand{\theosanswer}{\thechapter.\arabic{osanswer}}

\newenvironment{chapterexercise}[1]
{%
  \par\Needspace{4\baselineskip}%
  \refstepcounter{osexercise}%
  \par\addvspace{0.95em}%
  \noindent{\color{BookAccent}\bfseries 练习 \theosexercise}\quad\textbf{#1}\par\nobreak
  \smallskip\begingroup\leftskip=1.65em\rightskip=0pt\noindent\ignorespaces
}
{\par\endgroup}

\newenvironment{chapteranswer}[1]
{%
  \par\Needspace{4\baselineskip}%
  \refstepcounter{osanswer}%
  \par\addvspace{1.0em}%
  \noindent{\color{BookTeal}\bfseries 解答 \theosanswer}\quad\textbf{#1}\par\nobreak
  \smallskip\begingroup\leftskip=1.65em\rightskip=0pt\noindent\ignorespaces
}
{\par\endgroup}

\newenvironment{partbridge}
{%
  \clearpage
  \thispagestyle{plain}
  \vspace*{0.25em}
  {\noindent\Large\bfseries\color{BookNavy}本部分导读\par}%
  \vspace{0.55em}
}
{%
  \par\clearpage
}

\tikzset{
  os diagram/.style={
    font=\footnotesize\sffamily,
    text=BookInk,
    >=Latex,
    node distance=11mm and 14mm,
    line cap=round,
    line join=round
  },
  os block/.style={
    draw=BookRule,
    fill=white,
    rounded corners=1.5pt,
    line width=0.56pt,
    inner xsep=5.5pt,
    inner ysep=4.3pt,
    align=center,
    minimum height=8.5mm
  },
  os small block/.style={
    os block,
    font=\scriptsize\sffamily,
    inner xsep=4pt,
    inner ysep=3pt,
    minimum height=6.5mm
  },
  os bit/.style={
    os small block,
    minimum width=7mm,
    minimum height=6mm,
    inner sep=2pt
  },
  os state/.style={os block, draw=BookTeal, fill=BookGreenTint},
  os compute/.style={os block, draw=BookAccent, fill=BookBlueTint},
  os memory/.style={os block, draw=BookAmber, fill=BookAmberTint},
  os control/.style={os block, draw=BookRed, fill=BookRedTint},
  os bus/.style={
    draw=BookMuted,
    line width=0.58pt,
    shorten >=1.4pt,
    shorten <=1.4pt,
    -{Latex[length=2.25mm,width=1.85mm]}
  },
  os strong bus/.style={
    draw=BookAccent,
    line width=0.82pt,
    shorten >=1.4pt,
    shorten <=1.4pt,
    -{Latex[length=2.55mm,width=2.05mm]}
  },
  os feedback/.style={
    draw=BookRed,
    line width=0.72pt,
    shorten >=1.4pt,
    shorten <=1.4pt,
    -{Latex[length=2.4mm,width=1.95mm]},
    dashed
  },
  os guide/.style={draw=BookRule, line width=0.38pt},
  os label/.style={
    font=\scriptsize\sffamily,
    text=BookMuted,
    align=center,
    fill=white,
    inner sep=1.35pt
  },
  os title/.style={
    font=\scriptsize\bfseries\sffamily,
    text=BookNavy,
    align=left
  },
  os brace/.style={
    decorate,
    decoration={brace,amplitude=4pt},
    draw=BookAccent,
    line width=0.42pt
  },
  stage/.style={
    os compute,
    minimum width=3.5cm,
    minimum height=1.0cm
  },
  flow/.style={os strong bus}
}

\newenvironment{pdfdiagram}[1][]
{\par\addvspace{0.72em plus 0.18em}\goodbreak\begingroup\centering\begin{tikzpicture}[os diagram,scale=0.94,transform shape,#1]}
{\end{tikzpicture}\par\endgroup\nobreak\vspace{0.38em plus 0.10em}}

\newcommand{\diagramnote}[1]{%
  \par\nopagebreak[4]%
  {\noindent\footnotesize\color{BookMuted}图示：#1\par}%
  \addvspace{0.32em}%
}

\newcommand{\oscontract}[4]{%
  \begin{longtable}{|p{0.18\linewidth}|p{0.76\linewidth}|}
    \toprule
    状态 & #1 \\\tablerowrule
    入口 & #2 \\\tablerowrule
    失败 & #3 \\\tablerowrule
    证据 & #4 \\
    \bottomrule
  \end{longtable}
}

\newcommand{\tightlist}{%
  \setlength{\itemsep}{0.18em}%
  \setlength{\parsep}{0pt}%
  \setlength{\topsep}{0.2em}%
}

\newcommand{\detailsection}[1]{%
  \par\bigskip
  {\noindent\large\bfseries\color{BookNavy}#1\par}%
  \smallskip
}

\newcommand{\detailsubsection}[1]{%
  \par\smallskip
  {\noindent\bfseries\color{BookAccent}#1。\quad}%
  \ignorespaces
}

\newcommand{\detailsubsubsection}[1]{%
  \par\smallskip
  {\noindent\bfseries\color{BookMuted}#1。\quad}%
  \ignorespaces
}

% 演化样章使用行内主题标题连接相邻机器状态；正式卷同步这些确定性生成素材时
% 保留同一语义，避免把每个因果转折提升为新的编号小节。
\newcommand{\topic}[1]{%
  \par\Needspace{3\baselineskip}\medskip
  {\normalsize\bfseries\color{BookNavy}#1}\quad\ignorespaces
}

\newcommand{\inputtopic}[1]{%
  \begingroup
  \let\TopicSection\section
  \let\TopicSubsection\subsection
  \let\TopicSubsubsection\subsubsection
  \renewcommand{\chapter}[1]{\TopicSection{##1}}%
  \renewcommand{\section}[1]{\TopicSubsection{##1}}%
  \renewcommand{\subsection}[1]{\TopicSubsubsection{##1}}%
  \input{#1}%
  \endgroup
}

\newcommand{\inputdetail}[1]{%
  \begingroup
  \let\DetailSection\detailsection
  \let\DetailSubsection\detailsubsection
  \let\DetailSubsubsection\detailsubsubsection
  \renewcommand{\chapter}[1]{\DetailSection{##1}}%
  \renewcommand{\section}[1]{\DetailSection{##1}}%
  \renewcommand{\subsection}[1]{\DetailSubsection{##1}}%
  \renewcommand{\subsubsection}[1]{\DetailSubsubsection{##1}}%
  \input{#1}%
  \endgroup
}

% 把旧主章迁入新章节外壳：只去掉旧章标题，保留节级因果结构。
\newcommand{\inputchapterbody}[1]{%
  \begingroup
  \renewcommand{\chapter}[1]{}%
  \input{#1}%
  \endgroup
}

% 演进样本的每一章在正式卷第一章中成为一个局部演进阶段。
\newcommand{\inputsamplechapter}[1]{%
  \begingroup
  \let\SampleSection\section
  \let\SampleSubsection\subsection
  \let\SampleSubsubsection\subsubsection
  \renewcommand{\chapter}[1]{\SampleSection{##1}}%
  \renewcommand{\section}[1]{\SampleSubsection{##1}}%
  \renewcommand{\subsection}[1]{\SampleSubsubsection{##1}}%
  \input{#1}%
  \endgroup
}
